\documentclass[11pt,a4paper]{article}
\usepackage[utf8]{inputenc}
\usepackage[T1]{fontenc}
\usepackage[catalan]{babel}
\usepackage{lmodern}
\usepackage[sfdefault,lining]{FiraSans}
\usepackage[a4paper,top=2.3cm,bottom=2.2cm,left=1.8cm,right=1.8cm,headheight=1.3cm,headsep=0.4cm,footskip=1cm]{geometry}
\usepackage{amsmath,amssymb}
\usepackage{icomma}
\usepackage{graphicx}
\usepackage[table]{xcolor}
\usepackage{tikz}
\usepackage{fontawesome5}
\usepackage{ccicons}
\usepackage{fancyhdr}
\usepackage{lastpage}
\usepackage{tasks}
\usepackage{needspace}
\usepackage[most]{tcolorbox}
\usepackage{array}
\usepackage[hidelinks,pdftitle={Solucions de la fitxa de la unitat 5: Geometria analítica en el pla. 1r de Batxillerat MI},pdfsubject={Matemàtiques I},pdfkeywords={geometria analítica, vectors, producte escalar, equacions de la recta, posicions relatives, rectes perpendiculars, angles, distàncies, mediatriu, bisectriu}]{hyperref}

\definecolor{navy}{HTML}{1B2748}
\definecolor{accent}{HTML}{E8603A}
\definecolor{teal}{HTML}{1F9E90}
\definecolor{soft}{HTML}{EEF1F8}
\definecolor{muted}{HTML}{5D6685}

\setlength{\parindent}{0pt}
\setlength{\parskip}{3pt}
\settasks{label=\alph*), label-format=\bfseries\color{navy}, label-width=1.3em, item-indent=1.8em, column-sep=0.6em, after-item-skip=5pt, before-skip=3pt, after-skip=4pt}

% ---------- capçalera i peu ----------
\newcommand{\llicencia}{\href{https://creativecommons.org/licenses/by-nc-sa/4.0/deed.ca}{\ccbyncsa}\ \ Llicència Creative Commons Reconeixement-NoComercial-CompartirIgual 4.0 Internacional (CC BY-NC-SA 4.0)}
\pagestyle{fancy}
\fancyhf{}
\fancyhead[L]{\raisebox{-0.35\height}{\includegraphics[height=0.95cm]{logo-gva.png}}}
\fancyhead[C]{\small\color{muted}Solucions de la fitxa · Unitat 5: Geometria analítica en el pla · 1r Batxillerat MI}
\fancyhead[R]{\raisebox{-0.35\height}{\includegraphics[height=0.95cm]{logo-vermellar.png}}}
\fancyfoot[L]{\scriptsize\color{muted}\llicencia}
\fancyfoot[R]{\small\color{muted}\thepage}
\renewcommand{\headrulewidth}{0.4pt}
\renewcommand{\headrule}{\color{muted!40}\hrule height \headrulewidth}
\fancypagestyle{primera}{\fancyhf{}\renewcommand{\headrulewidth}{0pt}
  \fancyfoot[L]{\scriptsize\color{muted}\llicencia}
  \fancyfoot[R]{\small\color{muted}\thepage}}

% ---------- nivells i exercicis ----------
\newcommand{\estrelles}[1]{{\color{accent}\ifcase#1\or\faStar\faStar[regular]\faStar[regular]\or\faStar\faStar\faStar[regular]\or\faStar\faStar\faStar\fi}}
\newcounter{exer}
\newcommand{\ex}[1]{\par\needspace{4\baselineskip}\medskip\refstepcounter{exer}\noindent{\color{navy}\bfseries\theexer.}\ \ }

\newcommand{\seccio}[2]{\par\needspace{16\baselineskip}\bigskip
  \begin{tcolorbox}[colback=navy,colframe=navy,coltext=white,arc=2mm,boxrule=0pt,left=3mm,right=3mm,top=1.2mm,bottom=1.2mm]
  {\large\bfseries #1\quad #2}
  \end{tcolorbox}\vspace{-2pt}}

\newcommand{\exf}[1]{\par\medskip\refstepcounter{exer}\noindent{\color{navy}\bfseries\theexer.}\ \ }

% quadre de recordatori (només a la fitxa d'exercicis)
\newcommand{\recorda}[1]{\par\smallskip\begin{tcolorbox}[colback=accent!7,colframe=accent!55,arc=2mm,boxrule=0.5pt,left=2.5mm,right=2.5mm,top=1mm,bottom=1mm]\small{\color{accent}\faLightbulb[regular]}\ \textbf{Recorda:}\ #1\end{tcolorbox}}
\newcommand{\res}[1]{\boxed{#1}}\newcommand{\rest}[1]{\fbox{#1}}

\def\fmneed{11}% que el títol d'una secció no quede sol al final del full
% =====================================================================
% Capçalera FentMates de les fitxes: un color per a cada curs.
% Els colors són els de la pàgina inicial (LEVELS a assets/js/materials.js);
% el logo ve de _eines/logo/pdf/ (python3 _eines/logo/genera.py).
%
% Ús, al final del preàmbul de la unitat:
%   % =====================================================================
% Capçalera FentMates de les fitxes: un color per a cada curs.
% Els colors són els de la pàgina inicial (LEVELS a assets/js/materials.js);
% el logo ve de _eines/logo/pdf/ (python3 _eines/logo/genera.py).
%
% Ús, al final del preàmbul de la unitat:
%   % =====================================================================
% Capçalera FentMates de les fitxes: un color per a cada curs.
% Els colors són els de la pàgina inicial (LEVELS a assets/js/materials.js);
% el logo ve de _eines/logo/pdf/ (python3 _eines/logo/genera.py).
%
% Ús, al final del preàmbul de la unitat:
%   \input{../comu/fentmates.tex}
%   \fmcurs{1eso}
% i, al començament del document:
%   \capcalera{Matemàtiques}{1r ESO}{Fitxa d'exercicis · Unitat 2}{Divisibilitat}
% Opcional abans de \input: \def\fmneed{14} (espai mínim abans de cada secció).
% =====================================================================
\makeatletter
\def\fm@logodir{../../logo/pdf}
\providecommand{\fmneed}{10}
\def\fm@cursnom{}

% clau del curs, color i color del text damunt del color
\newcommand{\fm@curs}[3]{\@namedef{fm@color@#1}{#2}\@namedef{fm@text@#1}{#3}}
\fm@curs{1eso}{E6B84A}{navy}
\fm@curs{2eso}{3B82F6}{white}
\fm@curs{3eso}{FF7A50}{white}
\fm@curs{4eso-a}{A78BFA}{white}
\fm@curs{4eso-b}{6D28D9}{white}
\fm@curs{1batx-mi}{9B1C1C}{white}
\fm@curs{1batx-ccss}{EF6B6B}{white}
\fm@curs{2batx-mii}{1E6B47}{white}
\fm@curs{2batx-ccss}{2BB3A3}{white}

\newcommand{\fmcurs}[1]{%
  \@ifundefined{fm@color@#1}{\PackageError{fentmates}{Curs desconegut: #1}{}}{}%
  \edef\fm@tmp{\noexpand\definecolor{curs}{HTML}{\@nameuse{fm@color@#1}}%
               \noexpand\colorlet{curstext}{\@nameuse{fm@text@#1}}}\fm@tmp
  \colorlet{cursfosc}{curs!72!black}%
  \def\fm@logo{\fm@logodir/vertical-#1.pdf}%
  \def\fm@marca{\fm@logodir/marca-#1.pdf}}

% ---------- capçalera de la primera pàgina ----------
% #1 matèria (petit), #2 curs, #3 tipus de document i unitat, #4 títol
\newcommand{\capcalera}[4]{%
  \gdef\fm@cursnom{#2}%
  \par\noindent
  \begin{tcolorbox}[colback=curs,colframe=curs,boxrule=0pt,arc=2.2mm,width=3.5cm,height=2.15cm,
      valign=center,left=3mm,right=2mm,top=1mm,bottom=1mm,nobeforeafter,box align=center]
    \color{curstext}{\footnotesize #1}\par\vspace{1pt}{\fontsize{19}{21}\selectfont\bfseries #2\par}%
  \end{tcolorbox}\hspace{5mm}%
  \begin{minipage}[c]{\dimexpr\linewidth-3.5cm-5mm-2.9cm\relax}
    {\color{cursfosc}\small\bfseries\MakeUppercase{#3}}\par\vspace{2pt}
    {\color{navy}\fontsize{20}{23}\selectfont\bfseries #4\par}
  \end{minipage}\hfill
  \begin{minipage}[c]{2.6cm}\raggedleft\includegraphics[width=2.6cm]{\fm@logo}\end{minipage}\par
  \vspace{2pt}}

% ---------- barres de secció del color del curs ----------
\renewcommand{\seccio}[2]{\par\needspace{\fmneed\baselineskip}\bigskip
  \begin{tcolorbox}[colback=curs,colframe=curs,coltext=curstext,arc=2mm,boxrule=0pt,left=3mm,right=3mm,top=1.2mm,bottom=1.2mm]
  {\large\bfseries #1\quad #2}
  \end{tcolorbox}\vspace{-2pt}}

% ---------- capçalera i peu de les pàgines ----------
% primera pàgina: logos institucionals al peu
\fancypagestyle{primera}{\fancyhf{}\renewcommand{\headrulewidth}{0pt}
  \fancyfoot[L]{\raisebox{-0.32\height}{\includegraphics[height=0.78cm]{logo-gva.png}}\hspace{5mm}%
    \raisebox{-0.32\height}{\includegraphics[height=0.78cm]{logo-vermellar.png}}}
  \fancyfoot[R]{\scriptsize\color{muted}\href{https://creativecommons.org/licenses/by-nc-sa/4.0/deed.ca}{\ccbyncsa}\ \ CC BY-NC-SA 4.0 · fentmates.org\qquad\small\thepage}}
% la resta: el símbol a l'esquerra i el curs a la dreta, amb el títol al centre com abans
\fancyhead[L]{\raisebox{-0.3\height}{\includegraphics[height=0.62cm]{\fm@marca}}}
\fancyhead[R]{\small\bfseries\color{cursfosc}\fm@cursnom}
\renewcommand{\headrule}{\color{curs!45}\hrule height \headrulewidth}
\makeatother

%   \fmcurs{1eso}
% i, al començament del document:
%   \capcalera{Matemàtiques}{1r ESO}{Fitxa d'exercicis · Unitat 2}{Divisibilitat}
% Opcional abans de \input: \def\fmneed{14} (espai mínim abans de cada secció).
% =====================================================================
\makeatletter
\def\fm@logodir{../../logo/pdf}
\providecommand{\fmneed}{10}
\def\fm@cursnom{}

% clau del curs, color i color del text damunt del color
\newcommand{\fm@curs}[3]{\@namedef{fm@color@#1}{#2}\@namedef{fm@text@#1}{#3}}
\fm@curs{1eso}{E6B84A}{navy}
\fm@curs{2eso}{3B82F6}{white}
\fm@curs{3eso}{FF7A50}{white}
\fm@curs{4eso-a}{A78BFA}{white}
\fm@curs{4eso-b}{6D28D9}{white}
\fm@curs{1batx-mi}{9B1C1C}{white}
\fm@curs{1batx-ccss}{EF6B6B}{white}
\fm@curs{2batx-mii}{1E6B47}{white}
\fm@curs{2batx-ccss}{2BB3A3}{white}

\newcommand{\fmcurs}[1]{%
  \@ifundefined{fm@color@#1}{\PackageError{fentmates}{Curs desconegut: #1}{}}{}%
  \edef\fm@tmp{\noexpand\definecolor{curs}{HTML}{\@nameuse{fm@color@#1}}%
               \noexpand\colorlet{curstext}{\@nameuse{fm@text@#1}}}\fm@tmp
  \colorlet{cursfosc}{curs!72!black}%
  \def\fm@logo{\fm@logodir/vertical-#1.pdf}%
  \def\fm@marca{\fm@logodir/marca-#1.pdf}}

% ---------- capçalera de la primera pàgina ----------
% #1 matèria (petit), #2 curs, #3 tipus de document i unitat, #4 títol
\newcommand{\capcalera}[4]{%
  \gdef\fm@cursnom{#2}%
  \par\noindent
  \begin{tcolorbox}[colback=curs,colframe=curs,boxrule=0pt,arc=2.2mm,width=3.5cm,height=2.15cm,
      valign=center,left=3mm,right=2mm,top=1mm,bottom=1mm,nobeforeafter,box align=center]
    \color{curstext}{\footnotesize #1}\par\vspace{1pt}{\fontsize{19}{21}\selectfont\bfseries #2\par}%
  \end{tcolorbox}\hspace{5mm}%
  \begin{minipage}[c]{\dimexpr\linewidth-3.5cm-5mm-2.9cm\relax}
    {\color{cursfosc}\small\bfseries\MakeUppercase{#3}}\par\vspace{2pt}
    {\color{navy}\fontsize{20}{23}\selectfont\bfseries #4\par}
  \end{minipage}\hfill
  \begin{minipage}[c]{2.6cm}\raggedleft\includegraphics[width=2.6cm]{\fm@logo}\end{minipage}\par
  \vspace{2pt}}

% ---------- barres de secció del color del curs ----------
\renewcommand{\seccio}[2]{\par\needspace{\fmneed\baselineskip}\bigskip
  \begin{tcolorbox}[colback=curs,colframe=curs,coltext=curstext,arc=2mm,boxrule=0pt,left=3mm,right=3mm,top=1.2mm,bottom=1.2mm]
  {\large\bfseries #1\quad #2}
  \end{tcolorbox}\vspace{-2pt}}

% ---------- capçalera i peu de les pàgines ----------
% primera pàgina: logos institucionals al peu
\fancypagestyle{primera}{\fancyhf{}\renewcommand{\headrulewidth}{0pt}
  \fancyfoot[L]{\raisebox{-0.32\height}{\includegraphics[height=0.78cm]{logo-gva.png}}\hspace{5mm}%
    \raisebox{-0.32\height}{\includegraphics[height=0.78cm]{logo-vermellar.png}}}
  \fancyfoot[R]{\scriptsize\color{muted}\href{https://creativecommons.org/licenses/by-nc-sa/4.0/deed.ca}{\ccbyncsa}\ \ CC BY-NC-SA 4.0 · fentmates.org\qquad\small\thepage}}
% la resta: el símbol a l'esquerra i el curs a la dreta, amb el títol al centre com abans
\fancyhead[L]{\raisebox{-0.3\height}{\includegraphics[height=0.62cm]{\fm@marca}}}
\fancyhead[R]{\small\bfseries\color{cursfosc}\fm@cursnom}
\renewcommand{\headrule}{\color{curs!45}\hrule height \headrulewidth}
\makeatother

%   \fmcurs{1eso}
% i, al començament del document:
%   \capcalera{Matemàtiques}{1r ESO}{Fitxa d'exercicis · Unitat 2}{Divisibilitat}
% Opcional abans de \input: \def\fmneed{14} (espai mínim abans de cada secció).
% =====================================================================
\makeatletter
\def\fm@logodir{../../logo/pdf}
\providecommand{\fmneed}{10}
\def\fm@cursnom{}

% clau del curs, color i color del text damunt del color
\newcommand{\fm@curs}[3]{\@namedef{fm@color@#1}{#2}\@namedef{fm@text@#1}{#3}}
\fm@curs{1eso}{E6B84A}{navy}
\fm@curs{2eso}{3B82F6}{white}
\fm@curs{3eso}{FF7A50}{white}
\fm@curs{4eso-a}{A78BFA}{white}
\fm@curs{4eso-b}{6D28D9}{white}
\fm@curs{1batx-mi}{9B1C1C}{white}
\fm@curs{1batx-ccss}{EF6B6B}{white}
\fm@curs{2batx-mii}{1E6B47}{white}
\fm@curs{2batx-ccss}{2BB3A3}{white}

\newcommand{\fmcurs}[1]{%
  \@ifundefined{fm@color@#1}{\PackageError{fentmates}{Curs desconegut: #1}{}}{}%
  \edef\fm@tmp{\noexpand\definecolor{curs}{HTML}{\@nameuse{fm@color@#1}}%
               \noexpand\colorlet{curstext}{\@nameuse{fm@text@#1}}}\fm@tmp
  \colorlet{cursfosc}{curs!72!black}%
  \def\fm@logo{\fm@logodir/vertical-#1.pdf}%
  \def\fm@marca{\fm@logodir/marca-#1.pdf}}

% ---------- capçalera de la primera pàgina ----------
% #1 matèria (petit), #2 curs, #3 tipus de document i unitat, #4 títol
\newcommand{\capcalera}[4]{%
  \gdef\fm@cursnom{#2}%
  \par\noindent
  \begin{tcolorbox}[colback=curs,colframe=curs,boxrule=0pt,arc=2.2mm,width=3.5cm,height=2.15cm,
      valign=center,left=3mm,right=2mm,top=1mm,bottom=1mm,nobeforeafter,box align=center]
    \color{curstext}{\footnotesize #1}\par\vspace{1pt}{\fontsize{19}{21}\selectfont\bfseries #2\par}%
  \end{tcolorbox}\hspace{5mm}%
  \begin{minipage}[c]{\dimexpr\linewidth-3.5cm-5mm-2.9cm\relax}
    {\color{cursfosc}\small\bfseries\MakeUppercase{#3}}\par\vspace{2pt}
    {\color{navy}\fontsize{20}{23}\selectfont\bfseries #4\par}
  \end{minipage}\hfill
  \begin{minipage}[c]{2.6cm}\raggedleft\includegraphics[width=2.6cm]{\fm@logo}\end{minipage}\par
  \vspace{2pt}}

% ---------- barres de secció del color del curs ----------
\renewcommand{\seccio}[2]{\par\needspace{\fmneed\baselineskip}\bigskip
  \begin{tcolorbox}[colback=curs,colframe=curs,coltext=curstext,arc=2mm,boxrule=0pt,left=3mm,right=3mm,top=1.2mm,bottom=1.2mm]
  {\large\bfseries #1\quad #2}
  \end{tcolorbox}\vspace{-2pt}}

% ---------- capçalera i peu de les pàgines ----------
% primera pàgina: logos institucionals al peu
\fancypagestyle{primera}{\fancyhf{}\renewcommand{\headrulewidth}{0pt}
  \fancyfoot[L]{\raisebox{-0.32\height}{\includegraphics[height=0.78cm]{logo-gva.png}}\hspace{5mm}%
    \raisebox{-0.32\height}{\includegraphics[height=0.78cm]{logo-vermellar.png}}}
  \fancyfoot[R]{\scriptsize\color{muted}\href{https://creativecommons.org/licenses/by-nc-sa/4.0/deed.ca}{\ccbyncsa}\ \ CC BY-NC-SA 4.0 · fentmates.org\qquad\small\thepage}}
% la resta: el símbol a l'esquerra i el curs a la dreta, amb el títol al centre com abans
\fancyhead[L]{\raisebox{-0.3\height}{\includegraphics[height=0.62cm]{\fm@marca}}}
\fancyhead[R]{\small\bfseries\color{cursfosc}\fm@cursnom}
\renewcommand{\headrule}{\color{curs!45}\hrule height \headrulewidth}
\makeatother

\fmcurs{1batx-mi}

\begin{document}
\thispagestyle{primera}

\capcalera{Matemàtiques I}{1r Batx}{Solucions de la fitxa · Unitat 5}{\raggedright Geometria analítica en el pla}
\vspace{8pt}

\begin{tcolorbox}[colback=soft,colframe=soft,arc=2mm,left=3mm,right=3mm,top=1.5mm,bottom=1.5mm]
\small Intenta resoldre cada exercici abans de mirar-ne la solució. Els resultats es donen de manera exacta (fraccions i arrels) i, quan cal, amb el valor aproximat; els angles, en graus, minuts i segons. Pots comprovar-los amb els ginys dels apunts interactius de la unitat i amb el laboratori del triangle.
\end{tcolorbox}

% =====================================================
\seccio{1}{Vectors i operacions amb vectors}

\exf{1} \begin{tasks}(2)
  \task $\overrightarrow{AB}=\res{(-5,\,5)}$, $\overrightarrow{BA}=\res{(5,\,-5)}$
  \task $|\overrightarrow{AB}|=\sqrt{25+25}=\res{5\sqrt{2}}\approx 7{,}07$
  \task $\dfrac{(-5,\,5)}{5\sqrt2}=\res{\left(-\tfrac{\sqrt2}{2},\,\tfrac{\sqrt2}{2}\right)}$
  \task $10\cdot\left(\tfrac{\sqrt2}{2},-\tfrac{\sqrt2}{2}\right)=\res{(5\sqrt2,\,-5\sqrt2)}$
\end{tasks}

\ex{1} \begin{tasks}(3)
  \task $\res{45^\circ}$
  \task $\res{135^\circ}$
  \task $\res{225^\circ}$
  \task $\res{270^\circ}$
  \task $\res{0^\circ}$
  \task $\res{315^\circ}$
\end{tasks}La calculadora dona $\arctan\frac{v_2}{v_1}$ entre $-90^\circ$ i $90^\circ$: si $v_1<0$ cal sumar-hi $180^\circ$.

\ex{1} \begin{tasks}(3)
  \task $\res{(2\sqrt3,\,2)}$
  \task $\res{(-5\sqrt2,\,5\sqrt2)}$
  \task $\res{(-3,\,-3\sqrt3)}$
\end{tasks}

\ex{1} \begin{tasks}(2)
  \task $\res{(-4,\,1)}$
  \task $\res{(8,\,-17)}$
  \task $\res{(-\frac{1}{2},\,7)}$
  \task $\res{(9,\,14)}$
\end{tasks}

\ex{2} $x+2=5$ i $3+2y=-1$: $\res{x=3,\ y=-2}$. \quad $\sqrt{k^2+16}=5\Rightarrow k^2=9\Rightarrow\res{k=\pm3}$.

\ex{2} $\dfrac m6=\dfrac{-2}{3}\Rightarrow\res{m=-4}$; aleshores $(-4,\,-2)=\res{-\tfrac23}\,(6,\,3)$.

\ex{2} $\frac13\neq\frac{2}{-1}$: no són paral·lels, formen base. $a(1,\,2)+b(3,\,-1)=(1,\,9)\Rightarrow\left\{\begin{array}{l}a+3b=1\\2a-b=9\end{array}\right.\Rightarrow a=4,\ b=-1$: $\res{\vec w=4\vec u-\vec v}$.


% =====================================================
\seccio{2}{Aplicacions geomètriques dels vectors}

\exf{1} $M_{AB}=\res{(1,\,2)}$, \quad $M_{BC}=\res{(3,\,-1)}$, \quad $M_{CA}=\res{(0,\,-2)}$

\ex{1} $\frac{4+b_1}{2}=1$, $\frac{3+b_2}{2}=-2$: $\res{B(-2,\,-7)}$ (o bé $B=2M-A$).

\ex{1} $\overrightarrow{AB}=(12,\,-6)$, $\frac13\overrightarrow{AB}=(4,\,-2)$: $\res{M(1,\,2)}$, $\res{N(5,\,0)}$.

\ex{2} $\frac15\overrightarrow{AB}=(2,\,1)$: $\res{(3,\,0)},\ \res{(5,\,1)},\ \res{(7,\,2)},\ \res{(9,\,3)}$

\ex{1} \begin{tasks}(2)
  \task $\overrightarrow{AB}=(2,\,4)$, $\overrightarrow{AC}=(-3,\,-6)$: $\frac{-3}{2}=\frac{-6}{4}$, \rest{sí} estan alineats.
  \task $\overrightarrow{AB}=(3,\,2)$, $\overrightarrow{AC}=(6,\,5)$: $\frac63\neq\frac52$, \rest{no}.
\end{tasks}$\overrightarrow{AB}=(3,\,k-2)$, $\overrightarrow{AC}=(6,\,-6)$: $\frac{k-2}{3}=\frac{-6}{6}\Rightarrow k-2=-3\Rightarrow\res{k=-1}$.

\ex{1} $P'=2Q-P=\res{(-4,\,11)}$; \quad $R'=\res{(-5,\,-1)}$.

\ex{2} $\overrightarrow{AD}=\overrightarrow{BC}=(1,\,4)\Rightarrow\res{D(0,\,2)}$. El centre és el punt mig de les diagonals: $\res{(2,\,\frac{1}{2})}$.

\begin{center}\begin{tikzpicture}[scale=0.38,baseline=(current bounding box.center)]
\draw[muted!25,very thin] (-3,-4) grid (7,5);
\draw[-latex,muted] (-3.3,0)--(7.5,0) node[right,font=\scriptsize]{$x$}; \draw[-latex,muted] (0,-4.3)--(0,5.5) node[above,font=\scriptsize]{$y$};
\begin{scope}\clip (-3,-4) rectangle (7,5);
\filldraw[fill=teal!15,draw=teal,thick] (-1,-2)--(4,-1)--(5,3)--(0,2)--cycle;
\end{scope}
\fill[navy] (-1,-2) circle (2.2pt) node[below left,font=\scriptsize]{$A$};
\fill[navy] (4,-1) circle (2.2pt) node[below right,font=\scriptsize]{$B$};
\fill[navy] (5,3) circle (2.2pt) node[above right,font=\scriptsize]{$C$};
\fill[navy] (0,2) circle (2.2pt) node[above left,font=\scriptsize]{$D$};
\fill[navy] (2,0.5) circle (2.2pt) node[right,font=\scriptsize]{$M$};
\end{tikzpicture}\end{center}



% =====================================================
\seccio{3}{Producte escalar i angle de dos vectors}

\exf{1} $\vec u\cdot\vec v=12-2=\res{10}$; $|\vec u|=\res{\sqrt{13}}$, $|\vec v|=\res{\sqrt{17}}$; $\cos\alpha=\dfrac{10}{\sqrt{221}}\Rightarrow\res{\alpha\approx 47^\circ43'35''}$

\ex{1} \begin{tasks}(2)
  \task $\cos\alpha=\frac{0}{\sqrt{5}\cdot\sqrt{5}}\Rightarrow\res{90^\circ}$
  \task $\cos\alpha=\frac{-1}{\sqrt{2}\cdot1}\Rightarrow\res{135^\circ}$
  \task $\cos\alpha=\frac{-24}{5\cdot5}\Rightarrow\res{163^\circ44'23''}$
  \task $\cos\alpha=\frac{-1}{\sqrt{5}\cdot\sqrt{10}}\Rightarrow\res{98^\circ7'48''}$
\end{tasks}

\ex{2} \begin{tasks}(1)
  \task $2\cdot5\cdot\cos120^\circ=\res{-5}$
  \task $2|\vec u|^2-\vec u\cdot\vec v-|\vec v|^2=8+5-25=\res{-12}$
  \task $\sqrt{4+25+2(-5)}=\res{\sqrt{19}}$
\end{tasks}

\ex{2} $3m-12=0\Rightarrow\res{m=4}$. \quad $\cos60^\circ=\dfrac{1}{\sqrt{1+k^2}}=\dfrac12\Rightarrow1+k^2=4\Rightarrow\res{k=\pm\sqrt3}$.

\ex{1} Un perpendicular és $(8,\,6)$, de mòdul 10: $\res{\pm\left(\tfrac45,\tfrac35\right)}$.

\ex{2} $\overrightarrow{BA}\cdot\overrightarrow{BC}=(-4,\,-2)\cdot(-2,\,4)=0$: \rest{angle recte en $B$}. $|\overrightarrow{BA}|=|\overrightarrow{BC}|=\sqrt{20}$: isòsceles. Angles: $\res{90^\circ,\ 45^\circ,\ 45^\circ}$.

\begin{center}\begin{tikzpicture}[scale=0.38,baseline=(current bounding box.center)]
\draw[muted!25,very thin] (-1,-1) grid (7,8);
\draw[-latex,muted] (-1.3,0)--(7.5,0) node[right,font=\scriptsize]{$x$}; \draw[-latex,muted] (0,-1.3)--(0,8.5) node[above,font=\scriptsize]{$y$};
\begin{scope}\clip (-1,-1) rectangle (7,8);
\filldraw[fill=teal!15,draw=teal,thick] (1,1)--(5,3)--(3,7)--cycle;
\end{scope}
\fill[navy] (1,1) circle (2.2pt) node[below left,font=\scriptsize]{$A$};
\fill[navy] (5,3) circle (2.2pt) node[right,font=\scriptsize]{$B$};
\fill[navy] (3,7) circle (2.2pt) node[above,font=\scriptsize]{$C$};
\end{tikzpicture}\end{center}


\ex{2} $\dfrac{\vec u\cdot\vec v}{|\vec v|}=\dfrac{6-20}{5}=\res{-\tfrac{14}{5}}$: és negativa perquè l'angle és obtús.


% =====================================================
\seccio{4}{Equacions de la recta}

\exf{1} \begin{tasks}(1)
  \task Vectorial: $(x,\,y)=(2,\,-1)+\lambda(3,\,4)$
  \task Paramètriques: $\left\{\begin{array}{l}x=2+3\lambda\\ y=-1+4\lambda\end{array}\right.$
  \task Contínua: $\dfrac{x-2}{3}=\dfrac{y+1}{4}$
  \task Punt-pendent: $y+1=\frac{4}{3}(x-2)$
  \task Explícita: $y=\frac{4}{3}x-\frac{11}{3}$
  \task General: $\res{4x-3y-11=0}$
\end{tasks}

\ex{1} $\overrightarrow{AB}=(4,\,-8)\parallel(1,\,-2)$.\begin{tasks}(1)
  \task Vectorial: $(x,\,y)=(-1,\,3)+\lambda(1,\,-2)$
  \task Paramètriques: $\left\{\begin{array}{l}x=-1+\lambda\\ y=3-2\lambda\end{array}\right.$
  \task Contínua: $\dfrac{x+1}{1}=\dfrac{y-3}{-2}$
  \task Punt-pendent: $y-3=-2(x+1)$
  \task Explícita: $y=-2x+1$
  \task General: $\res{2x+y-1=0}$
\end{tasks}

\ex{1} $y+2=-\tfrac32(x-4)$; \quad $y=-\tfrac32x+4$; \quad $\res{3x+2y-8=0}$

\ex{1} $m=\res{\tfrac25}$; talls: $\res{(-5,\,0)}$ i $\res{(0,\,2)}$; $\vec u=\res{(5,\,2)}$, $\vec n=\res{(2,\,-5)}$; $\left\{\begin{array}{l}x=5\lambda\\ y=2+2\lambda\end{array}\right.$

\ex{1} $P$: $\lambda=2$ en les dues equacions: \rest{$P\in r$}. $Q$: $\lambda=1$ i $\lambda=\tfrac32$: \rest{$Q\notin r$}. $5=-3+2\lambda\Rightarrow\lambda=4\Rightarrow\res{k=13}$.

\ex{1} \begin{tasks}(2)
  \task $\res{y=-2}$
  \task $\res{x=3}$
  \task $\res{2x+3y=0}$
  \task $\res{y=x}$
\end{tasks}

\ex{1} $3(x-1)-(y-5)=0\Rightarrow\res{3x-y+2=0}$

\ex{2} $m=\tan135^\circ=-1$: $y-3=-(x-2)\Rightarrow\res{x+y-5=0}$. \quad $m=\tan60^\circ=\sqrt3$: $\res{y=\sqrt3\,x-4}$.


% =====================================================
\seccio{5}{Posicions relatives de dues rectes}

\exf{1} \begin{tasks}(1)
  \task $\frac24\neq\frac{-3}{1}$: \rest{secants} en $\res{(1,\,1)}$.
  \task $\frac{1}{-3}=\frac{-2}{6}=\frac{3}{-9}$: \rest{coincidents}.
  \task $\frac36=\frac12\neq\frac{-2}{1}$: \rest{paral·leles}.
\end{tasks}

\ex{1} Directors $(2,\,1)$ i $(-4,\,2)$, no proporcionals: \rest{secants}. $\left\{\begin{array}{l}1+2\lambda=3-4\mu\\ -1+\lambda=2\mu\end{array}\right.\Rightarrow\mu=0,\ \lambda=1$: punt de tall $\res{(3,\,0)}$.

\ex{1} $\vec u_r\cdot\vec n_s=(1,\,-3)\cdot(3,\,1)=0$: paral·leles o coincidents. $(2,\,1)\in r$ i $6+1-7=0$: \rest{coincidents}.

\ex{2} a) $\frac a2\neq\frac{3}{-1}$: $\res{a\neq-6}$. b) $\res{a=-6}$ (i $\frac{-2}{5}\neq-3$). \rest{No}: amb $a=-6$ els termes independents no són proporcionals.

\ex{2} $\frac m3=\frac{2}{-1}=\frac n4\Rightarrow\res{m=-6,\ n=-8}$.

\ex{2} $r\cap s=\res{(2,\,3)}$, \quad $r\cap t=\res{(-2,\,-1)}$, \quad $s\cap t=\res{(6,\,-1)}$

\begin{center}\begin{tikzpicture}[scale=0.38,baseline=(current bounding box.center)]
\draw[muted!25,very thin] (-4,-3) grid (8,6);
\draw[-latex,muted] (-4.3,0)--(8.5,0) node[right,font=\scriptsize]{$x$}; \draw[-latex,muted] (0,-3.3)--(0,6.5) node[above,font=\scriptsize]{$y$};
\begin{scope}\clip (-4,-3) rectangle (8,6);
\draw[navy,thick] (-4,-3)--(8,9);
\draw[accent,thick] (-4,9)--(8,-3);
\draw[teal,thick] (-4,-1)--(8,-1);
\end{scope}
\fill[navy] (2,3) circle (2.2pt) node[above,font=\scriptsize]{$$};
\fill[navy] (-2,-1) circle (2.2pt) node[below,font=\scriptsize]{$$};
\fill[navy] (6,-1) circle (2.2pt) node[below,font=\scriptsize]{$$};
\end{tikzpicture}\end{center}



% =====================================================
\seccio{6}{Paral·lelisme i perpendicularitat}

\exf{1} Paral·lela: $3x-2y+K=0$, $6+2+K=0$: $\res{3x-2y-8=0}$. \quad Perpendicular: $2x+3y+K=0$, $4-3+K=0$: $\res{2x+3y-1=0}$.

\ex{1} Paral·lela (director $(-2,\,1)$): $\res{x+2y+4=0}$. \quad Perpendicular (director $(1,\,2)$): $\res{2x-y-2=0}$.

\ex{1} Paral·lela ($m=-\tfrac13$): $\res{y=-\tfrac13x+3}$. \quad Perpendicular ($m'=3$): $\res{y=3x+13}$.

\ex{2} $(k,\,-2)\cdot(3,\,k+1)=3k-2k-2=0\Rightarrow\res{k=2}$. Paral·leles: $\frac k3=\frac{-2}{k+1}\Rightarrow k^2+k+6=0$, sense solució real: \rest{mai no ho són}.

\ex{2} Altura: perpendicular a $\overrightarrow{AB}=(6,\,2)$ per $C$: $\res{3x+y-16=0}$. Recta $AB$: $x-3y+2=0$. Peu: $\res{H(\frac{23}{5},\,\frac{11}{5})}$.

\ex{2} $M(1,\,-1)$, $\vec n=\overrightarrow{AB}=(8,\,-6)\parallel(4,\,-3)$: $\res{4x-3y-7=0}$. \par Definició: $(x+3)^2+(y-2)^2=(x-5)^2+(y+4)^2\Rightarrow16x-12y-28=0$ \checkmark


% =====================================================
\seccio{7}{Angle de dues rectes}

\exf{1} $\cos\alpha=\dfrac{|(3,\,-1)\cdot(1,\,2)|}{\sqrt{10}\sqrt5}=\dfrac{1}{\sqrt{50}}\Rightarrow\res{\alpha\approx 81^\circ52'12''}$

\ex{1} $\cos\alpha=\dfrac{|(2,\,-1)\cdot(-1,\,3)|}{\sqrt5\sqrt{10}}=\dfrac{5}{\sqrt{50}}=\dfrac{\sqrt2}{2}\Rightarrow\res{\alpha=45^\circ}$

\ex{1} Directors $(1,\,2)$ i $(1,\,-1)$: $\cos\alpha=\dfrac{1}{\sqrt{10}}\Rightarrow\res{\alpha\approx 71^\circ33'54''}$

\ex{1} a) $m=1$: $\res{45^\circ}$. \quad b) $m=-\tfrac12$: la recta forma $\res{26^\circ33'54''}$ amb l'eix (l'angle d'inclinació és $153^\circ26'6''$).

\ex{3} $m_r=\tfrac13$. $\left|\dfrac{m-\frac13}{1+\frac m3}\right|=1\Rightarrow m=2$ o $m=-\tfrac12$: $\res{2x-y=0}$ i $\res{x+2y-5=0}$.


% =====================================================
\seccio{8}{Distàncies}

\exf{1} $|AB|=\res{5}$, $|BC|=\res{5}$, $|CA|=\res{5\sqrt2}$: isòsceles i, com que $25+25=50$, \rest{rectangle} en $B$. Perímetre: $\res{10+5\sqrt2}\approx17{,}07$.

\ex{1} \begin{tasks}(2)
  \task $\dfrac{|12+6+2|}{5}=\res{4}$
  \task $2x-y+1=0$: $\dfrac{9}{\sqrt5}=\res{\tfrac{9\sqrt5}{5}}\approx4{,}02$
  \task $\res{4}$
  \task $x+y-3=0$: $\dfrac{|-2|}{\sqrt2}=\res{\sqrt2}$
\end{tasks}

\ex{1} $\frac36=\frac48\neq\frac{-5}{15}$: paral·leles. $(-1,\,2)\in r$: $d=\dfrac{|-6+16+15|}{10}=\res{\tfrac52}$.

\ex{2} $(0,\,y)$: $\dfrac{|-4y+1|}{5}=2\Rightarrow-4y+1=\pm10$: $\res{\left(0,-\tfrac94\right)}$ i $\res{\left(0,\tfrac{11}{4}\right)}$.

\ex{2} $5x-12y+K=0$ amb $\dfrac{|K-3|}{13}=2$: $\res{5x-12y+29=0}$ i $\res{5x-12y-23=0}$.

\ex{2} El radi és la distància del centre a la recta: $\dfrac{|6-4-12|}{5}=\res{2}$.


% =====================================================
\seccio{9}{Problemes mètrics}

\exf{2} Perpendicular per $P$: $(x,\,y)=(-3,\,4)+\lambda(2,\,-1)$. En $r$: $5\lambda-5=0\Rightarrow\lambda=1$: $\res{M(-1,\,3)}$, \ $P'=2M-P=\res{(1,\,2)}$.

\begin{center}\begin{tikzpicture}[scale=0.38,baseline=(current bounding box.center)]
\draw[muted!25,very thin] (-5,-1) grid (4,6);
\draw[-latex,muted] (-5.3,0)--(4.5,0) node[right,font=\scriptsize]{$x$}; \draw[-latex,muted] (0,-1.3)--(0,6.5) node[above,font=\scriptsize]{$y$};
\begin{scope}\clip (-5,-1) rectangle (4,6);
\draw[navy,thick] (-5,-5)--(4,13);
\draw[accent,thick] (-3,4)--(1,2);
\end{scope}
\fill[navy] (-3,4) circle (2.2pt) node[left,font=\scriptsize]{$P$};
\fill[navy] (-1,3) circle (2.2pt) node[below right,font=\scriptsize]{$M$};
\fill[navy] (1,2) circle (2.2pt) node[right,font=\scriptsize]{$P'$};
\end{tikzpicture}\end{center}


\ex{1} Respecte de $y=x$: $\res{(1,\,4)}$ (s'intercanvien les coordenades). Respecte de $x=2$: $\res{(0,\,1)}$.

\ex{3} $\dfrac{3x+4y-1}{5}=\pm\dfrac{5x-12y+2}{13}$: $\res{14x+112y-23=0}$ i $\res{64x-8y-3=0}$. \ $(14,\,112)\cdot(64,\,-8)=0$ \checkmark

\ex{2} Vèrtexs $(2,\,3)$, $(-2,\,-1)$, $(6,\,-1)$. Base sobre $t$: $8$; altura: $d\big((2,\,3),t\big)=4$. $S=\dfrac{8\cdot4}{2}=\res{16\ \text{u}^2}$.

\ex{2} \begin{tasks}(2)
  \task $M_{BC}=\left(\tfrac52,\,4\right)$: $\res{8x-9y+16=0}$
  \task $\perp\overrightarrow{BC}=(-3,\,4)$: $\res{3x-4y+6=0}$
  \task $\res{G(1,\,\frac{8}{3})}$
  \task $\tfrac12|6\cdot6-2\cdot3|=\res{15\ \text{u}^2}$
\end{tasks}

\ex{2} Mediatrius de $AB$ i $AC$: $x=3$, $y=2$: $\res{O(3,\,2)}$; $R=d(O,\,A)=\res{\sqrt{13}}$. (És el punt mig de la hipotenusa $BC$: el triangle és rectangle.)

\ex{2} Altura de $C$: $x=1$. Altura de $A$ ($\perp\overrightarrow{BC}=(-2,\,4)$): $x-2y+3=0$. Ortocentre: $\res{H(1,\,2)}$.

\ex{2} Centre $(3,\,3)$; la mitja diagonal $(2,\,2)$ girada $90^\circ$ és $(2,\,-2)$: $\res{B(5,\,1)}$, $\res{D(1,\,5)}$. Àrea: $\dfrac{d^2}{2}=\dfrac{32}{2}=\res{16\ \text{u}^2}$.

\begin{center}\begin{tikzpicture}[scale=0.38,baseline=(current bounding box.center)]
\draw[muted!25,very thin] (-1,-1) grid (7,7);
\draw[-latex,muted] (-1.3,0)--(7.5,0) node[right,font=\scriptsize]{$x$}; \draw[-latex,muted] (0,-1.3)--(0,7.5) node[above,font=\scriptsize]{$y$};
\begin{scope}\clip (-1,-1) rectangle (7,7);
\filldraw[fill=teal!15,draw=teal,thick] (1,1)--(5,1)--(5,5)--(1,5)--cycle;
\draw[navy,thick] (1,1)--(5,5);
\end{scope}
\fill[navy] (1,1) circle (2.2pt) node[below left,font=\scriptsize]{$A$};
\fill[navy] (5,1) circle (2.2pt) node[below right,font=\scriptsize]{$B$};
\fill[navy] (5,5) circle (2.2pt) node[above right,font=\scriptsize]{$C$};
\fill[navy] (1,5) circle (2.2pt) node[above left,font=\scriptsize]{$D$};
\end{tikzpicture}\end{center}


\ex{3} És a la mediatriu de $AB$: $M(3,\,1)$, $\vec n=(4,\,2)\parallel(2,\,1)$: $2x+y-7=0$. Amb $y=x$: $\res{(\frac{7}{3},\,\frac{7}{3})}$.


\end{document}
